\section{Reading the Evidence Map}\label{sec:taxonomy}
The same ticket can expose different facts at the network link, the tool recipient and the model service.
\subsection{From private task to observed execution}
An execution maps a hidden task or state through an agent policy and environment into communication, then into an observer's features:
\begin{center}
\begin{tabular}{c}
hidden task/state $\longrightarrow$ agent execution\\
$\longrightarrow$ observed communication $\longrightarrow$ inference.
\end{tabular}
\end{center}
Different cuts through this chain expose different objects. A client link can reveal response mode or hosted-agent identity. Tool egress exposes contacted services and request/response shape. The model link carries prompts, tool feedback and model responses. Inter-agent links expose delegation and coordination. Figure~\ref{fig:chain} places these observations on the feedback path.




\subsection{What the inferred label means}
The ticket example separates four labels. \emph{Intent} is the requested objective or assigned task family. \emph{Execution} is the recorded action sequence. \emph{Success} states whether that sequence and its final state meet the task's criterion. \emph{Authorization} states whether the action is permitted for the principal under the applicable policy. A failed refund task can still contain a refund attempt; a successful refund can still be unauthorized. An evaluation should bind its target to the corresponding reference instead of letting one label stand for all four.

Exposure is relative to both the target and the observer. The same execution may expose intended-task information on M and actual contacted services on T. Initial instructions, tool feedback and later decisions all contribute to the complete M trace. Assigning a task-family label to that trace measures intention recognition; attributing the signal to a particular feedback step requires a further controlled comparison.


\subsection{What each record specifies}
Six fields give a score its meaning: the \emph{interface and visible representation}; \emph{observer access}; \emph{prior information} supplied by the evaluator; \emph{target}; \emph{population and base rate}; and \emph{protected training access}. These fields define the comparison record. The matrix records the fields once; the reporting checklist in Section~\ref{sec:reporting} turns them into entries for future studies.

Access determines what a feature can mean. A server can inspect decrypted HTTP/2 properties that a transit observer cannot. The content-aware MCP study's ``metadata'' includes argument hashes, and its graph uses content matches~\cite[Secs.~4.1--4.2]{zavrak2026content}. Its representation therefore requires more than packet metadata. Logs and browser instrumentation can provide evaluation labels without becoming attacker inputs.

ANT's supplied stage boundary is prior information; its 2,101 stages and 47 labels define the decision population; the primitive is the target~\cite{li2026ant}. In the live OHTTP comparison, an unchanged random-forest (RF) family refits on protected data, so it is an adapted attacker~\cite{dangol2026workflow}.

\subsection{Deployment determines who can see each signal}
Table~\ref{tab:deployment-access} maps common deployments to observation rights and the corresponding source evidence. A visible model link is an access condition, not a property granted to every attacker. Captured packets, TLS records, HTTP streams, server-sent events (SSE) and tool calls are different units. Segmentation, buffering and retransmission change their correspondence. Supplement S7 gives the protocol-level map.

Two simultaneous tickets may use separate connections or share an encrypted HTTP/2 connection. In the latter case, a transit observer sees combined packets while stream identifiers remain encrypted. A gateway that buffers SSE output changes the pattern again. The connection and flush policy therefore determine the measured trace.

\begin{table*}[t]\centering\small
\caption{Deployment and access map. The configurations are analyst-defined; cited settings test the listed view, not every component of the deployment. A gateway's plaintext access requires TLS termination.}\label{tab:deployment-access}
\setlength{\tabcolsep}{4pt}
\begin{tabularx}{\textwidth}{@{}P{.18\textwidth}P{.25\textwidth}P{.28\textwidth}Y@{}}\toprule
Deployment & Observer and permitted view & Existing evidence for that view & Protection consequence \\\midrule
Enterprise egress & A transit monitor sees routed model/tool flows, sizes, times and any exposed destination association. & WRA: supplied research-domain traces; ANT and cloud benchmark: model/combined traffic~\cite{jeong2026wra,li2026ant,pouryousef2026traffic}. & Specify which services cross the monitored egress. A union of flows is available only when the monitor can observe them. \\\addlinespace[3pt]
Shared gateway or relay & Transit sees multiplexed gateway output. A terminating gateway can additionally read addressed requests. & Live OHTTP: passive-link inference under stated relay trust~\cite{dangol2026workflow}; PrivacyTrace: semantic recipient/gateway access~\cite{duan2026privacytrace}. & Separate relay association protection from content trust. Connection sharing changes the observation unit. \\\addlinespace[3pt]
Local model, remote tools & An external observer sees remote tool traffic; the local model call has no corresponding external network leg. & WRA's destination-recovery mechanism is relevant to tool egress~\cite{jeong2026wra}; whole-deployment protection is untested in the mapped comparison. & Evaluate tool and downstream calls. Do not import a remote-model attack score into an absent network leg. \\\addlinespace[3pt]
Cloud model, remote tools & An authorized vantage may see either or both encrypted legs. A contacted tool sees its own arguments. & ANT/cloud benchmark test model traffic~\cite{li2026ant,pouryousef2026traffic}; PrivacyTrace tests recipient-visible representations~\cite{duan2026privacytrace}. & Check permitted model and tool views after the intervention; recipients and passive monitors need different protections. \\\bottomrule
\end{tabularx}
\end{table*}


The same reasoning explains defense placement. Padding a tool message before a gateway does not specify the gateway's output. Local MCP over \texttt{stdio} removes that network tool leg, while remote model requests may remain. Shared key/value (KV) cache state creates an additional timing interface for a querying tenant. The map therefore records the transport, observation layer, terminating parties and shared state alongside each result. Section~\ref{sec:live-study} compares real HTTP/2 concurrency with paired feedback and shaping interventions.
